\chapter{أطوار القمر}\label{ch:g7:moon-phases}

منذ سنوات مضت كنت تمسك دفتر \omterm{def:g2:moon-in-the-sky:moon}{قمر} وتراقب الشكل الفضّي يزحف من
\omterm{ex:g2:moon-in-the-sky:shapes}{الهلال} إلى \omterm{ex:g2:moon-in-the-sky:shapes}{البدر} ثم يعود، ووُعدت بجواب ``لماذا'' في وقت لاحق.
واللاحق هو الآن. فعندك كل ما يحتاج إليه التفسير: \omterm{def:g2:moon-in-the-sky:moon}{قمر} يدور
شهريًّا، \omterm{def:g1:light-and-shadows:source}{وضوء} شمس في خطوط مستقيمة، وكرة نصفها مضاء دائمًا ---
دائمًا.

\section{الحقيقة الواحدة تحت كل شيء}

\begin{proposition}[نصفه مضاء، دائمًا]\label{prop:g7:moon-phases:halflit}
تضيء الشمس نصف \omterm{def:g2:moon-in-the-sky:moon}{القمر} في كل لحظة --- النصف المواجه لها ---
تمامًا كما تضيء نصف الأرض. فليس \omterm{def:g2:moon-in-the-sky:moon}{للقمر} شكل متغيّر، ولا نموّ ولا
انكماش: هو كرة حجرية تلبس جانب نهار دائمًا وجانب ليل دائمًا،
والحدّ بينهما ينزلق وهو يدور. وكل ما نسمّيه ``طورًا'' ليس إلا
\emph{المنظر}: كم من جانب النهار يصادف أن يواجه الأرض.
\end{proposition}

\begin{definition}[الأطوار]\label{def:g7:moon-phases:phases}
\emph{أطوار}\index{طور} \omterm{def:g2:moon-in-the-sky:moon}{القمر} هي أشكاله الظاهرة على \omterm{def:g6:sun-earth-moon:axisorbit}{مدار} الشهر:
\emph{المحاق} (جانب النهار مُدار عنّا كله --- فلا نرى شيئًا)،
و\emph{\omterm{ex:g2:moon-in-the-sky:shapes}{الهلال}}، و\emph{التربيع الأول} (نصف القرص مضاء)،
و\emph{الأحدب} (أكثر من النصف)، و\emph{\omterm{ex:g2:moon-in-the-sky:shapes}{البدر}} (جانب النهار
مُدار نحونا كله)، ثم الأحدب، و\emph{التربيع الأخير} \omterm{ex:g2:moon-in-the-sky:shapes}{والهلال}
ثانية، فتنغلق الدورة في نحو $29.5$ يومًا. ومن المحاق إلى \omterm{ex:g2:moon-in-the-sky:shapes}{البدر}
يكبر الجزء المضاء ليلة بعد ليلة: \omterm{def:g2:moon-in-the-sky:moon}{فالقمر} \emph{متزايد}\index{تزايد}؛
ومن \omterm{ex:g2:moon-in-the-sky:shapes}{البدر} رجوعًا إلى المحاق يصغر: \emph{متناقص}\index{تناقص}.
\end{definition}

\begin{omfigure}
\begin{tikzpicture}[scale=0.8]
  % dashed orbit
  \draw[dashed, omExo!70] (0,0) ellipse (4.2 and 3.1);
  % eight stations: dark half toward the left, sunlit half (white)
  % toward the sunlight arriving from the right
  \foreach \a in {0,45,...,315} {
    \begin{scope}[shift={({4.2*cos(\a)},{3.1*sin(\a)})}]
      \fill[omExo] (0,0) circle (0.46);
      \begin{scope}
        \clip (0,0) circle (0.46);
        \fill[white] (0,-0.5) rectangle (0.5,0.5);
      \end{scope}
      \draw[thick] (0,0) circle (0.46);
      \draw[thick] (0,0.46) -- (0,-0.46);
    \end{scope}
  }
  % the Earth, its own day side facing the sunlight
  \fill[omDef!35] (0,0) circle (0.55);
  \begin{scope}
    \clip (0,0) circle (0.55);
    \fill[omDef!12] (0,-0.6) rectangle (0.6,0.6);
  \end{scope}
  \draw[thick] (0,0) circle (0.55);
  \node[font=\footnotesize] at (0,0) {الأرض};
  \foreach \y in {-2.4,-1.2,0,1.2,2.4} {
    \draw[-Stealth, omThm, thick] (7.7,\y) -- (6.4,\y);
  }
  \node[right, font=\small, omThm] at (7.8,0.55) {ضوء الشمس};
  \node[left, font=\footnotesize] at (-4.85,0) {بدر};
  \node[right, font=\footnotesize] at (4.85,0) {محاق};
  \node[above, font=\footnotesize] at (0,3.75) {التربيع الأخير};
  \node[below, font=\footnotesize] at (0,-3.75) {التربيع الأول};
\end{tikzpicture}

{\small السرّ كله في صورة واحدة: ثماني محطّات \omterm{def:g2:moon-in-the-sky:moon}{للقمر} الدائر، كل
واحدة نصفها المواجه للشمس مضاء. والمتغيّر الوحيد أيّ شريحة من
جانب النهار تواجه الأرض.}
\end{omfigure}

\begin{method}[كن الأرض، واحمل القمر]\label{met:g7:moon-phases:walk}
التفسير، مُؤدّى في دقيقة واحدة:
\begin{enumerate}
  \item في غرفة مظلمة، \omterm{def:g3:first-electric-circuit:bulb}{مصباح} عارٍ واحد هو الشمس؛ وأنت الأرض؛
        وكرة بيضاء على يدك الممدودة هي \omterm{def:g2:moon-in-the-sky:moon}{القمر}؛
  \item واجه \omterm{def:g3:first-electric-circuit:bulb}{المصباح} والكرة نحوه: فنصفها المضاء مُدار عنك ---
        وترى جانبها المظلم: محاق؛
  \item در ببطء في مكانك (والكرة تدور حولك): فتظهر شظيّة مضاءة
        عند حافة الكرة --- \omterm{ex:g2:moon-in-the-sky:shapes}{هلال} \omterm{def:g7:moon-phases:phases}{متزايد}؛ وبعد ربع دورة يظهر نصف
        الكرة مضاءً: التربيع الأول؛
  \item واصل الدوران: فالأحدب --- ثم، \omterm{def:g3:first-electric-circuit:bulb}{والمصباح} خلف ظهرك، الكرة
        مضاءة كلها أمامك: بدر؛ وأتمم الدورة عبر الأحدب \omterm{def:g7:moon-phases:phases}{المتناقص}
        \omterm{ex:g2:moon-in-the-sky:shapes}{والهلال} إلى المحاق.
\end{enumerate}
دورة واحدة منك تعيد تمثيل شهر واحد من السماء.
\end{method}

\section{الوهم الشائع، متقاعدًا}

\begin{proposition}[الأطوار ليست ظلّ الأرض]\label{prop:g7:moon-phases:notshadow}
أشهر جواب خاطئ في علم الفلك: ``\omterm{def:g7:moon-phases:phases}{الأطوار} هي \omterm{def:g1:light-and-shadows:shadow}{ظلّ} الأرض على
\omterm{def:g2:moon-in-the-sky:moon}{القمر}''. والهندسة تدحضه من ثلاث جهات. \omterm{def:g1:light-and-shadows:shadow}{فظلّ} الأرض يشير
\emph{بعيدًا} عن الشمس --- ولا يستطيع أن يلمس إلا قمرًا واقفًا
مقابل الشمس، وهذا لا يقع إلا عند \omterm{ex:g2:moon-in-the-sky:shapes}{البدر} (وحينها في الخسوفات
النادرة وحدها). \omterm{def:g2:moon-in-the-sky:moon}{والقمر} الهلالي يقف في جهة الشمس منّا، لا قرب
ظلّنا البتة. وحدّ \omterm{def:g1:light-and-shadows:shadow}{الظلّ} على \omterm{def:g2:moon-in-the-sky:moon}{القمر} جزء من دائرة كبيرة دائمًا ---
فما كان له أن ينحت حدّ الشكل الأحدب المزدوج الانحناء. \omterm{def:g7:moon-phases:phases}{والأطوار}
لا تحتاج إلى \omterm{def:g1:light-and-shadows:shadow}{ظلّ}: فكرة نصفها مضاء، منظورًا إليها من جهات
مختلفة، تفعل كل شيء.
\end{proposition}

\begin{example}[الخسوف مقابل الطور]\label{ex:g7:moon-phases:eclipsevs}
احفظ إعتامَي \omterm{def:g2:moon-in-the-sky:moon}{القمر} في درجين منفصلين. \emph{\omterm{def:g7:moon-phases:phases}{الطور}}: الدورة
الشهرية البطيئة، جانب \omterm{def:g2:moon-in-the-sky:moon}{القمر} الليلي نفسه يستدير نحونا --- ولا
دخل \omterm{def:g1:light-and-shadows:shadow}{لظلّ} لنا فيه. و\emph{\omterm{def:g7:shadows-eclipses:lunar}{الخسوف}}: ساعة نادرة يعبر فيها \omterm{ex:g2:moon-in-the-sky:shapes}{البدر}
ظلّنا التامّ --- \omterm{def:g1:light-and-shadows:shadow}{ظلّ} الأرض واقعًا عليه حقًّا، نحاسيًّا سريعًا.
فإيقاع الدفتر منظر؛ \omterm{def:g7:shadows-eclipses:lunar}{والخسوف} زيارة.
\end{example}

\section{قراءة السماء كالساعة}

\begin{example}[أين تحفظ الأطوار مواقيتها]\label{ex:g7:moon-phases:hours}
المخطط يثبّت كل \omterm{def:g7:moon-phases:phases}{طور} في موضع --- وكل موضع في جدول مواقيت. \omterm{ex:g2:moon-in-the-sky:shapes}{فالبدر}
يقف مقابل الشمس: فلا بدّ أن يطلع عند \omterm{def:g1:day-and-night:daynight}{الغروب}، ويملك الليل كله،
ويغيب عند \omterm{def:g1:day-and-night:daynight}{الشروق}. والمحاق يسافر \emph{مع} الشمس: طالعًا النهار
كله، ضائعًا في الوهج --- ولهذا لا ``يرى'' أحد محاقًا. والتربيع
الأول يتخلّف عن الشمس ربع دورة: فيطلع قرب منتصف النهار ويغيب
قرب منتصف الليل --- \omterm{def:g2:moon-in-the-sky:moon}{قمر} ما بعد العشاء. \omterm{def:g2:moon-in-the-sky:moon}{وقمر} النهار الذي فاجأك
في طفولتك صار الآن جدول مواقيت: فكل \omterm{def:g2:moon-in-the-sky:moon}{قمر} ليس محاقًا ولا بدرًا
يقاسم الشمس جزءًا من نوبتها.
\end{example}

\begin{example}[في أيّ جهة يشير الهلال؟]\label{ex:g7:moon-phases:points}
الجانب المضاء من أيّ \omterm{def:g7:moon-phases:phases}{طور} يواجه الشمس --- \omterm{ex:g2:moon-in-the-sky:shapes}{والهلال} قوس مصوَّب نحو
راميه. \omterm{ex:g2:moon-in-the-sky:shapes}{فهلال} المساء في الغرب: قرناه يشيران إلى أعلى وبعيدًا عن
توهّج \omterm{def:g1:day-and-night:daynight}{الغروب} تحت الأفق. والسماء لا ترسمه خطأً أبدًا؛ والرسّامون
يفعلون كثيرًا. \omterm{ex:g2:moon-in-the-sky:shapes}{وهلال} قرناه نحو الشمس هو لعبة ``جِد الخطأ''
المفضّلة عند علم الفلك.
\end{example}

\begin{remark}[الوجه الذي لا يستدير أبدًا]\label{rem:g7:moon-phases:sameface}
ملاحظة أخرى من الدفتر، تُوفّى حقّها أخيرًا: بقع \omterm{def:g2:moon-in-the-sky:moon}{القمر} الداكنة
ترسم الوجه نفسه كل ليلة، طورًا بعد \omterm{def:g7:moon-phases:phases}{طور}، وسنة بعد سنة. \omterm{def:g2:moon-in-the-sky:moon}{فالقمر}
يدور حول محوره دورة واحدة بالضبط في كل دورة مدارية --- فالنصف
نفسه يواجه الأرض إلى الأبد. ولم يرَ أحد الجانب البعيد حتى طار
مسبار حول \omterm{def:g2:moon-in-the-sky:moon}{القمر} بآلة تصوير، في ذاكرة جدّيك الحيّة. وأمّا لماذا
يتّفق الدوران \omterm{def:g6:sun-earth-moon:axisorbit}{والمدار} هذا الاتّفاق التامّ فقصة عميقة عن المدّ
والجزر، مؤجّلة إلى سنوات الجامعة من هذا المنهج.
\end{remark}

\section{تمارين}

\begin{exercise}[$\star$]\label{exo:g7:moon-phases:1}
ما الجزء من \omterm{def:g2:moon-in-the-sky:moon}{القمر} المضاء \omterm{def:g1:light-and-shadows:source}{بضوء} الشمس في أيّ لحظة؟ وما ``\omterm{def:g7:moon-phases:phases}{الطور}''
إذن؟
\end{exercise}

\begin{exercise}[$\star$]\label{exo:g7:moon-phases:2}
اسرد \omterm{def:g7:moon-phases:phases}{الأطوار} بالترتيب ابتداءً من المحاق، وأعطِ طول الدورة. وأيّ
نصفَي الدورة هو ``\omterm{def:g7:moon-phases:phases}{التزايد}''؟
\end{exercise}

\begin{exercise}[$\star$]\label{exo:g7:moon-phases:3}
في طريقة المشي المداري، أين \omterm{def:g3:first-electric-circuit:bulb}{المصباح} والكرة وأنت عند المحاق؟
وعند \omterm{ex:g2:moon-in-the-sky:shapes}{البدر}؟
\end{exercise}

\begin{exercise}[$\star$]\label{exo:g7:moon-phases:4}
أعطِ اثنين من الدحوض الهندسية الثلاثة لقول ``\omterm{def:g7:moon-phases:phases}{الأطوار} هي \omterm{def:g1:light-and-shadows:shadow}{ظلّ}
الأرض''.
\end{exercise}

\begin{exercise}[$\star$]\label{exo:g7:moon-phases:5}
ما الفرق بين ظلمة المحاق وظلمة \omterm{def:g7:shadows-eclipses:lunar}{الخسوف}؟
\end{exercise}

\begin{exercise}[$\star$]\label{exo:g7:moon-phases:6}
لماذا يطلع \omterm{ex:g2:moon-in-the-sky:shapes}{البدر} عند \omterm{def:g1:day-and-night:daynight}{الغروب} دائمًا؟ استدلّ من موضعه في المخطط.
\end{exercise}

\begin{exercise}[$\star$]\label{exo:g7:moon-phases:7}
لماذا لا ``يرى'' أحد المحاق يطلع في منتصف الليل؟ وأين المحاق في
منتصف الليل؟
\end{exercise}

\begin{exercise}[$\star$]\label{exo:g7:moon-phases:8}
\omterm{ex:g2:moon-in-the-sky:shapes}{هلال} مساء معلّق في الغرب. في أيّ جهة يشير قرناه بالنسبة إلى
الشمس الغائبة --- ولماذا لا يستطيع \omterm{ex:g2:moon-in-the-sky:shapes}{هلال} أن يصوّب قرنيه إلى
الشمس أبدًا؟
\end{exercise}

\begin{exercise}[$\star\star$]\label{exo:g7:moon-phases:9}
\omterm{def:g2:moon-in-the-sky:moon}{القمر} الليلة أحدب \omterm{def:g7:moon-phases:phases}{متزايد}. أين يقف في مداره؟ ومتى طلع تقريبًا،
وأيّ \omterm{def:g7:moon-phases:phases}{طور} --- وأيّ جدول مواقيت --- يأتي بعد خمس ليالٍ؟
\end{exercise}

\begin{exercise}[$\star\star$]\label{exo:g7:moon-phases:10}
لا يقع \omterm{def:g7:shadows-eclipses:solar}{الكسوف} إلا عند \omterm{def:g7:moon-phases:phases}{طور} واحد، ولا يقع \omterm{def:g7:shadows-eclipses:lunar}{الخسوف} إلا عند \omterm{def:g7:moon-phases:phases}{طور}
آخر. سمّهما، وبرّر بهندسة الاصطفاف، واشرح لماذا تبقى الخسوفات
والكسوفات نادرة عند هذين الطورين رغم ذلك.
\end{exercise}

\begin{exercise}[$\star\star$]\label{exo:g7:moon-phases:11}
روّاد فضاء على الجانب القريب من \omterm{def:g2:moon-in-the-sky:moon}{القمر} ينظرون إلى الأرض فوقهم.
صِف ``\omterm{def:g7:moon-phases:phases}{أطوار} الأرض'' التي يرونها على \omterm{def:g6:sun-earth-moon:axisorbit}{مدار} شهر، وأيّ \omterm{def:g7:moon-phases:phases}{طور} تُريهم
الأرض حين نرى نحن محاقًا. (أجرِ الهندسة في الاتجاهين.)
\end{exercise}

\begin{exercise}[$\star\star\star$]\label{exo:g7:moon-phases:12}
سرّ \omterm{ex:g2:moon-in-the-sky:shapes}{الهلال} القديم: في أمسيات صافية يتوهّج الجزء \emph{المظلم}
من \omterm{def:g2:moon-in-the-sky:moon}{القمر} الهلالي توهّجًا خافتًا --- ``\omterm{def:g2:moon-in-the-sky:moon}{القمر} القديم في حضن \omterm{def:g2:moon-in-the-sky:moon}{القمر}
الجديد''. لقيَ نفسُك في الصف الثاني هذا لغزًا؛ فحُلَّه الآن
كاملًا: تتبّع \omterm{def:g6:describing-motion:trajectory}{مسار} \omterm{def:g1:light-and-shadows:source}{الضوء}، واشرح لماذا يبلغ \omterm{def:g1:light-and-shadows:source}{ضوء} الأرض أشدّه عند
\omterm{def:g7:moon-phases:phases}{الأطوار} الهلالية بالضبط (أيّ \omterm{def:g7:moon-phases:phases}{طور} تُري \emph{الأرضُ} القمرَ
حينها؟)، ولماذا يخبو نحو \omterm{ex:g2:moon-in-the-sky:shapes}{البدر}.
\end{exercise}

\section{مسألة: شهر مع القمر}

\begin{problem}[{مسألة نهاية الأسبوع --- شهر المرصد المفتوح؛
ثلاثون ليلة من القمر، ولوحة جدارية واحدة؛ اختبار
المرشد}]\label{pb:g7:moon-phases:1}
يقيم مرصد البلدة برنامجًا عامًّا اسمه ``شهر مع \omterm{def:g2:moon-in-the-sky:moon}{القمر}''. وأنت،
بصفتك مرشدًا صغيرًا، تمسك اللوحة الجدارية وتجيب الزوّار.

\medskip\noindent\textbf{القسم الأول --- إعداد اللوحة.}
\begin{enumerate}
  \item الليلة $0$ محاق. املأ محطّات اللوحة: في أيّ الليالي
        تقريبًا يقع التربيع الأول \omterm{ex:g2:moon-in-the-sky:shapes}{والبدر} والتربيع الأخير؟
  \item وإلى جانب كل محطّة تطلب اللوحة أوقات الطلوع. طابِق:
        يطلع عند \omterm{def:g1:day-and-night:daynight}{الغروب}؛ \ يطلع قرب منتصف النهار؛ \ يطلع قرب
        منتصف الليل؛ \ يطلع ويغيب مع الشمس.
  \item يسأل زائر لماذا تُظهر اللوحة \omterm{ex:g2:moon-in-the-sky:shapes}{الهلال} منخفضًا في سماء
        المساء \emph{الغربية} بينما \omterm{ex:g2:moon-in-the-sky:shapes}{البدر} يتسلّق من
        \emph{الشرق} عند الغسق. اشرح بمواضع \omterm{def:g6:sun-earth-moon:axisorbit}{المدار}.
  \item ويقسم زائر آخر أنه رأى \omterm{def:g2:moon-in-the-sky:moon}{القمر} في العاشرة صباحًا. أيّ
        \omterm{def:g7:moon-phases:phases}{أطوار} اللوحة تجعل \omterm{def:g2:moon-in-the-sky:moon}{قمر} الصباح ممكنًا؟
\end{enumerate}

\medskip\noindent\textbf{القسم الثاني --- عروض المرشد.}
\begin{enumerate}[resume]
  \item تؤدّي مشي \omterm{def:g3:first-electric-circuit:bulb}{المصباح} والكرة لمجموعة مدرسية. فيوقفك تلميذ
        في منتصف الدورة: الكرة تُظهر ثلاثة أرباعها مضاءة. سمِّ
        \omterm{def:g7:moon-phases:phases}{الطور}، وحدّد للمجموعة مواضع \omterm{def:g3:first-electric-circuit:bulb}{المصباح} وأنت الأرض والكرة.
  \item ``إذن \omterm{def:g1:light-and-shadows:shadow}{ظلّ} الأرض هو الذي يصنع \omterm{def:g7:moon-phases:phases}{الأطوار}؟'' يسأل التلميذ
        مشيرًا إلى كرتك ذات الجانب المظلم. أعطِ دحض العرض
        نفسه: جانب مَن الليلي يرونه، وأين يكون \omterm{def:g1:light-and-shadows:shadow}{ظلّ} الأرض
        فعلًا؟
  \item وتسأل المجموعة لماذا يرون دائمًا ``الأرنب'' (أو
        ``الوجه'') في \omterm{def:g2:moon-in-the-sky:moon}{القمر}، ولا يرون ظهره أبدًا. أجب بحقيقة
        الوجه نفسه.
  \item ويقول أب متشكّك: ``إن كان \omterm{def:g2:moon-in-the-sky:moon}{القمر} يدور حول نفسه كما
        تزعم، أفلا ينبغي أن نراه يستدير؟'' حُلّ التناقض
        الظاهر: كيف يستطيع \omterm{def:g2:moon-in-the-sky:moon}{القمر} أن يدور دورة واحدة بالضبط في
        كل دورة مدارية فيُخفي بذلك دورانه عنّا؟
\end{enumerate}

\medskip\noindent\textbf{القسم الثالث --- ليلة الاختبار.}
\begin{enumerate}[resume]
  \item \omterm{def:g5:everyday-energy:energy}{بطاقة} الاختبار أ: ``\omterm{def:g2:moon-in-the-sky:moon}{قمر} الليلة طلع في منتصف الليل
        ويُظهر نصفه \emph{الأيسر} مضاءً (كما يُرى مع اقتراب
        الصباح). \omterm{def:g7:moon-phases:phases}{متزايد} أم \omterm{def:g7:moon-phases:phases}{متناقص} --- وأيّ تربيع؟''
  \item \omterm{def:g5:everyday-energy:energy}{بطاقة} الاختبار ب: ``ملصق فنّان يُظهر قمرًا هلاليًّا في
        منتصف الليل، عاليًا فوق الرؤوس، وقرناه يشيران مستقيمًا
        إلى أسفل نحو ألوان \omterm{def:g1:day-and-night:daynight}{الغروب} على الأفق.'' اسرد كل خطأ
        فلكي في الملصق.
  \item \omterm{def:g5:everyday-energy:energy}{بطاقة} الاختبار ج: ``أثناء بدر الليلة، هل يمكن أن يشهد
        الزوّار كسوفًا في مكان ما على الأرض أيضًا؟'' احسم
        الأمر بهندسة \omterm{def:g7:moon-phases:phases}{الأطوار}.
  \item كلمة الختام: في ثلاث جمل، أعطِ الحقيقة الواحدة الكامنة
        وراء الشهر، والوهم الذي دفنه الشهر، وحيلة جدول
        المواقيت التي ينبغي أن يحملها الزوّار إلى بيوتهم (ماذا
        يقول أيّ \omterm{def:g7:moon-phases:phases}{طور} عن ساعة طلوعه).
\end{enumerate}
\end{problem}
